 %%%%%%%%%%%%%%%%%%%%%%%%%%%%%%%%%%%%%%%%%%%%%%%%%%%%%%%%%%%%%%%%%%%%%%%%%%%%%%%%
%2345678901234567890123456789012345678901234567890123456789012345678901234567890
%        1         2         3         4         5         6         7         8

\documentclass[a4paper, 10pt, conference]{ieeeconf}

\IEEEoverridecommandlockouts                              % This command is only needed if
 % you want to use the \thanks command

\overrideIEEEmargins                                      % Needed to meet printer requirements.

%In case you encounter the following error:
%Error 1010 The PDF file may be corrupt (unable to open PDF file) OR
%Error 1000 An error occurred while parsing a contents stream. Unable to analyze the PDF file.
%This is a known problem with pdfLaTeX conversion filter. The file cannot be opened with acrobat reader
%Please use one of the alternatives below to circumvent this error by uncommenting one or the other
%\pdfobjcompresslevel=0
%\pdfminorversion=4


% The following packages can be found on http:\\www.ctan.org
\usepackage[utf8]{inputenc}
\usepackage[T1]{fontenc}
\usepackage[ngerman]{babel}
\usepackage{graphicx} % for pdf, bitmapped graphics files
%\usepackage{epsfig} % for postscript graphics files
%\usepackage{mathptmx} % assumes new font selection scheme installed
%\usepackage{times} % assumes new font selection scheme installed
\usepackage{amsmath} % assumes amsmath package installed
%\usepackage{amssymb}  % assumes amsmath package installed

%%%%%%%%%%%%% added for project course control engineering
\usepackage[left=1.57cm,right=1.57cm,top=2.54cm,bottom=2.54cm]{geometry}
\usepackage{subfigure}
%%%%%%%%%%%%% added for project course control engineering


\usepackage{amssymb}
\let\labelindent\relax % ieeeconf.cls already defines \labelindent (line 1444),
                       % which makes enumitem stop with "Command \labelindent
                       % already defined"; enumitem re-creates it itself.
\usepackage{enumitem}
\usepackage{multirow}
\let\proof\relax\let\endproof\relax % ieeeconf.cls already defines a proof
                       % environment -> amsthm stops with "Command \proof already
                       % defined"; drop the class version, amsthm provides its own.
\usepackage{amsthm}

%%%%%%%%%%%%% added for the Thermal Digital Twin (DT4TM) report %%%%%%%%%%%%%
\usepackage{booktabs}
\usepackage{tikz}
\usetikzlibrary{arrows.meta, positioning, fit, backgrounds}
\usepackage{pgfplots} % Abb. C1/C2 drawn from embedded data -- no image files needed
\pgfplotsset{compat=1.17}
\usepackage{url} % \url{} breaks long filenames at _ . - / so they wrap inside narrow table columns (plain \texttt{} cannot break there)
\usepackage{tabularx} % X columns share \textwidth exactly and auto-wrap text,
                       % unlike plain tabular's c/l columns which size to content
                       % and silently overflow the page margin if content is wide
\usepackage[hidelinks]{hyperref} % makes \cite/\ref/\eqref clickable (jump to the
                       % bibliography/label); loaded last so it patches every
                       % preceding package's cross-referencing commands. `hidelinks`
                       % keeps links invisible (no colored text/boxes) so the printed
                       % look is unchanged -- only PDF click-navigation is added.
% ieeeconf.cls defines a FAKE \NAT@parse purely to make hyperref think natbib
% is loaded (so hyperref's natbib-compatibility code, not its plain-\cite code,
% handles \cite) -- but the class's own comment says this fools hyperref into
% NOT hyperlinking citation numbers at all. The class explicitly names this
% override as the fix (ieeeconf.cls, \NAT@parse definition): undefine it so
% \cite{...} actually links down to the matching \bibitem.
\makeatletter
\let\NAT@parse\undefined
\makeatother
% Short helpers used in the text (plain LaTeX, no extra packages):
% \py{...}  -> "Python:" pointer under each equation (where it is solved)
% \anh{...} -> bold run-in heading for the numbered items A.1, B.2, ... in the appendices
% \src{...} -> line ID in report_numbers_output.txt (e.g. [R5])
\newcommand{\py}[1]{\par{\raggedright\noindent$\rightarrow$~\textit{Python:}~#1\par}\smallskip}
\newcommand{\anh}[1]{\par\medskip\noindent\textbf{#1}\quad}
\newcommand{\src}[1]{\,[#1]}
\newcolumntype{Y}{>{\raggedright\arraybackslash}X}
%%%%%%%%%%%%% added for the Thermal Digital Twin (DT4TM) report %%%%%%%%%%%%%

\theoremstyle{remark}
% Template default is \newtheorem{remark}{Remark}; renamed to the German
% "Anmerkung" only because this report is written in German. Same environment,
% same style -- no format change.
\newtheorem{remark}{Anmerkung}

% ieeeconf.cls hard-codes the English word "Appendix" for auto-generated
% appendix headings (see \@IEEEprocessthesectionargument); since this report
% is in German, override it here to "Anhang" so headings match the prose.
\makeatletter
\def\@IEEEprocessthesectionargument#1{%
\@ifmtarg{#1}{%
\@IEEEappendixsavesection*{Anhang \thesectiondis}%
\addcontentsline{toc}{section}{Anhang \thesection}}{%
\@IEEEappendixsavesection*{Anhang \thesectiondis \\* #1}%
\addcontentsline{toc}{section}{Anhang \thesection: #1}}}
\makeatother

\title{\LARGE \bf
Digital Twin für das Thermomanagement eines elektrodynamischen Levitators
}
\author{ Dang Minh Hoang, Marina Borchert, Mohamed Aziz El Majid % <-this % stops a space
\thanks{Project Course Digital Twin, Summer Term 2026. Supervisors: Hr. Prof. Dr. rer. nat. Dirk Hartmann, Fr. Dr.-Ing. Melina Merkel. Stand: 27.09.2026, Version~3.}% project course: only change the supervisor(s)
}


\begin{document}



\maketitle
\thispagestyle{plain}
\pagestyle{plain}


%%%%%%%%%%%%%%%%%%%%%%%%%%%%%%%%%%%%%%%%%%%%%%%%%%%%%%%%%%%%%%%%%%%%%%%%%%%%%%%%
\begin{abstract}
Ziel ist ein Modell, das die Temperatur von Aluminiumscheibe, Spulen und
Eisen eines TEAM-28-ähnlichen Levitators aus dem Spulenstrom berechnet, und
zwar schnell genug für eine laufende Darstellung. Die Physik wird als
Feldgleichungen formuliert und mit der Finite-Elemente-Methode (FEM)
gelöst. Dieses genaue, aber langsame Modell wird anschließend durch
Modellordnungsreduktion (engl.\ \emph{Model Order Reduction}, MOR) in ein
reduziertes Modell (engl.\ \emph{Reduced-Order Model}, ROM) mit acht
Zustandsgrößen überführt. Das
Ergebnis läuft als Python-Simulation und wird als interaktive HTML-Seite im
Browser dargestellt. Im stationären Zustand trifft das Modell die gemessenen
Spulentemperaturen, auf die es kalibriert ist. Den zeitlichen Verlauf trifft
es noch nicht; die Ursachen werden benannt.
\end{abstract}

%%%%%%%%%%%%%%%%%%%%%%%%%%%%%%%%%%%%%%%%%%%%%%%%%%%%%%%%%%%%%%%%%%%%%%%%%%%%%%%


\section{Einleitung}
\label{sec:intro}

Der Versuchsstand ist ein elektrodynamischer Levitator nach dem
TEAM-Benchmark~28~\cite{team28}. Zwei Spulen mit 50-Hz-Wechselstrom erzeugen
ein Magnetfeld. Es induziert in einer Aluminiumscheibe Ströme, die die
Scheibe anheben. Dieselben Ströme erzeugen Wärme in Scheibe, Spulen und
Eisen.

Diese Temperaturen sind schwer zu messen: Ein Fühler an der schwebenden
Scheibe würde die Levitation stören, und das Innere der Wicklung ist nicht
zugänglich. Der Spulenstrom dagegen ist leicht messbar. Ein Modell, das aus
dem Strom die Temperatur berechnet, wirkt daher als \textbf{virtueller
Sensor}~\cite{hartmann2025golden}. Es muss drei Anforderungen erfüllen:
\begin{enumerate}[nosep]
  \item \textbf{schnell:} Ergebnisse in Echtzeit, ohne Wartezeit;
  \item \textbf{physikalisch begründet:} aus den Grundgleichungen
        abgeleitet, nicht nur an Messkurven angepasst;
  \item \textbf{kalibrierbar:} wenige unsichere Parameter, die sich an
        Messungen nachstellen lassen.
\end{enumerate}
Genaue Modelle sind aber groß und langsam. Das Bindeglied ist die
\textbf{Modellordnungsreduktion} (MOR)~\cite{hartmann2018mor}: Sie
verkleinert das genaue Modell zu einem \textbf{reduzierten Modell} (ROM),
das in Echtzeit läuft. MOR ist also das Verfahren, das ROM sein Ergebnis.
Dieser Bericht zeigt beides am konkreten Beispiel.

\textbf{Aufbau.} Abschnitt~\ref{sec:problem} beschreibt das Problem,
Abschnitt~\ref{sec:method} die Methodik (FEM und Reduktion),
Abschnitt~\ref{sec:code} den Code, Abschnitt~\ref{sec:results} die
Ergebnisse, Abschnitt~\ref{sec:outlook} den Ausblick. Zusatzrechnungen und
Code-Details stehen in den Anhängen.


%%%%%%%%%%%%%%%%%%%%%%%%%%%%%%%%%%%%%%%%%%%%%%%%%%%%%%%%%%%%%%%%%%%%%%%%%%%%%%%
\section{Problemstellung}
\label{sec:problem}

Gesucht ist die \textbf{Temperaturverteilung im Gerät und ihr zeitlicher
Verlauf}, wenn ein Wechselstrom durch die Spulen fließt. Dazu wird das Gerät
zuerst \textbf{physikalisch} beschrieben: Welche Bauteile gibt es, wo
entsteht Wärme, wie wird sie abgegeben (Abschnitt~\ref{sec:phys})? Danach
wird jeder physikalische Vorgang in eine \textbf{Gleichung} übersetzt
(Abschnitt~\ref{sec:math}). Diese Gleichungen löst
Abschnitt~\ref{sec:method}.

\subsection{Physikalisches Modell}
\label{sec:phys}

\textbf{Aufbau.} Von innen nach außen: Eisenkern, Innenspule, Eisenring,
Außenspule (Abbildung~\ref{fig:geometry}). Darüber schwebt die
Aluminiumscheibe. Die Anordnung ist drehsymmetrisch um die senkrechte Achse.

\begin{figure}[t]
\centering
\begin{tikzpicture}[x=0.065cm, y=0.065cm,
    ironfill/.style={fill=gray!45},
    cufill/.style={fill=orange!35},
    alfill/.style={fill=cyan!25},
    lbl/.style={font=\scriptsize, align=center}]
  % Symmetrieachse
  \draw[dash dot] (0,15) -- (0,-58);
  \node[font=\tiny, right] at (2,16) {$r=0$};
  % r-Achse
  \draw[-{Stealth[length=4pt]}] (0,-58) -- (109,-58) node[below, font=\tiny] {$r$ [mm]};
  \draw (25.9,-58) -- (25.9,-56.5); \node[font=\tiny] at (25.9,-63) {25{,}9};
  \draw (61.9,-58) -- (61.9,-56.5); \node[font=\tiny] at (61.9,-63) {61{,}9};
  \draw (79.9,-58) -- (79.9,-56.5); \node[font=\tiny] at (79.9,-63) {79{,}9};
  \draw (102.9,-58) -- (102.9,-56.5); \node[font=\tiny] at (102.9,-63) {102{,}9};
  % Eisenkern
  \draw[ironfill] (0,0) rectangle (25.9,-53);
  \node[lbl] at (13,-26.5) {Eisen-\\kern};
  % Innenspule
  \draw[cufill] (27.9,0) rectangle (61.9,-53);
  \node[lbl] at (44.9,-26.5) {Innenspule\\ca.\ 1000 Wdg.\\Kupfer};
  % Eisenring
  \draw[ironfill] (64.9,0) rectangle (79.9,-53);
  \node[lbl, rotate=90] at (72.4,-26.5) {Eisenring};
  % Außenspule
  \draw[cufill] (82.9,0) rectangle (102.9,-53);
  \node[lbl] at (92.9,-26.5) {Außen-\\spule\\ca.\ 500\\Wdg.};
  % Al-Scheibe (schematische Position)
  \draw[alfill] (0,6) rectangle (80,9);
  \node[lbl] at (40,12.5) {Al-Scheibe, $R{=}80$, $d{=}3$};
  \draw[<->, dashed] (85,0) -- (85,6);
  \node[font=\tiny, right] at (86,3) {Luftspalt};
  % Höhe
  \draw[<->] (106,0) -- (106,-53);
  \node[font=\tiny, rotate=90] at (109,-26.5) {53};
\end{tikzpicture}
\caption{Halber Querschnitt (Maße in mm, aus \texttt{params.yaml}).}
\label{fig:geometry}
\end{figure}

\textbf{Wirkprinzip.} Der Wechselstrom erzeugt ein Magnetfeld, das sich
50-mal pro Sekunde umpolt. Das Eisen bündelt das Feld. Das wechselnde Feld
induziert in Scheibe und Eisen Kreisströme (Wirbelströme). Alle Ströme
erzeugen Wärme: in den Spulen der Spulenstrom, in Scheibe und Eisen die
Wirbelströme. Die Wärme verteilt sich im Material und geht an der Oberfläche
an die Luft.

\textbf{Vereinfachungen.}
\begin{itemize}[nosep]
  \item \textbf{Drehsymmetrie:} Es genügt, eine Halbebene $(r,z)$ zu rechnen.
  \item \textbf{Konstante Umgebung:} $20\,^\circ$C, wie mit den Betreuern
        vereinbart.
  \item \textbf{Konstanter Strom} als Standardfall ($5\,\mathrm A$). Im
        Modell ist der Strom frei einstellbar.
  \item \textbf{Feste Scheibenhöhe} in der Feldrechnung (Folgen:
        Abschnitt~\ref{sec:causes}).
\end{itemize}

\textbf{Eingangsdaten.} Tabelle~\ref{tab:inputs} zeigt die wichtigsten
Eingangswerte und woher sie stammen. Die vollständige Liste steht in
Anhang~A.1.

\begin{table}[t]
\centering\footnotesize
\caption{Eingangsdaten des Modells}
\label{tab:inputs}
\begin{tabularx}{\columnwidth}{@{}Y>{\raggedright\arraybackslash}p{2.3cm}Y@{}}
\toprule
Größe & Wert & Herkunft \\
\midrule
Radien von Kern, Spulen, Ring & 0--102{,}9\,mm (Abb.~\ref{fig:geometry}) & gemessen (Lineal) \\
Höhe Kern/Spulen/Ring & 53\,mm & gemessen \\
Windungszahl innen / außen & ca.\ 1000 / 500 & Angabe zum Aufbau \\
Scheibe: Radius / Dicke / Masse & 80\,mm / 3\,mm / 159--163\,g & gemessen \\
Spulenstrom (Effektivwert) & 5{,}0\,A (Betrieb), 7{,}8\,A (Kalibrierung) & gemessen (Multimeter PeakTech 2005) \\
Frequenz & 50\,Hz & Netz \\
Leitfähigkeit, Wärmedaten Al/Cu & Standardwerte & Literatur \\
Permeabilität Eisen $\mu_r$ & 1000 & \textbf{angenommen} \\
Wärmeübergänge an Scheibe und Luft & Schätzwerte & \textbf{angenommen} \\
Wärmeübergang der Spulen $hA$ & 2{,}47 / 2{,}89\,W/K & \textbf{kalibriert} (Abschnitt~\ref{sec:calib}) \\
Spulentemperatur bei 7{,}8\,A (innen / außen) & 79 / 74\,$^\circ$C & gemessen (IR-Kamera, einmalig) \\
\bottomrule
\end{tabularx}
\end{table}

Das Gerät hat \textbf{keine fest eingebauten Temperatursensoren}. Alle
Temperaturmesswerte stammen aus einzelnen Aufnahmen mit einer IR-Kamera
(23.06.2026).

\subsection{Mathematisches Modell}
\label{sec:math}

Jeder Vorgang aus Abschnitt~\ref{sec:phys} wird durch eine Gleichung
beschrieben. Tabelle~\ref{tab:model} am Ende fasst sie zusammen.

\subsubsection{Magnetfeld}
Berechnet wird das \textbf{magnetische Vektorpotential} $A_\varphi(r,z)$, aus
dem sich das Feld ableitet. Da der Strom sinusförmig ist, genügen Amplitude
und Phase, zusammengefasst in einer komplexen Zahl, dem
\textbf{Phasor}~\cite{biro1989vectorpotential}:
\begin{equation}
-\nabla\cdot(\nu\nabla A_\varphi) + \frac{\nu}{r^2}A_\varphi + j\omega\sigma A_\varphi = J_s
\label{eq:em}
\end{equation}
mit $\nu=1/(\mu_r\mu_0)$ und $\omega = 2\pi f$. \emph{In Worten:} Term~1
beschreibt die Ausbreitung des Felds (im Eisen besonders leicht). Term~2
folgt aus den Zylinderkoordinaten. Term~3 ist die Gegenwirkung der
Wirbelströme. Rechts steht die Quelle, die Stromdichte in den Spulen.
\py{\url{em_solver.py} (Anhang~B.2)}

\subsubsection{Wärmequelle}
Aus dem Feld folgt die Wirbelstromdichte $J = -j\omega\sigma A_\varphi$ und
daraus die mittlere Wärmeleistung pro Volumen:
\begin{equation}
q(r,z) = \tfrac12\,\sigma\,\omega^2\,|A_\varphi|^2 \quad[\mathrm{W/m^3}].
\label{eq:loss}
\end{equation}
In den Spulen gilt $P_{\mathrm{Cu}}=\tfrac12 I^2R$.
\py{\url{em_solver.py}, Funktion \url{compute_losses} (Anhang~B.2)}

\subsubsection{Wärmeleitung}
Die Wärme verteilt sich im Material und wird an der Oberfläche abgegeben:
\begin{equation}
\begin{gathered}
\rho c_p\frac{\partial T}{\partial t}=\nabla\cdot(k\nabla T)+q,\\
-k\frac{\partial T}{\partial n}=h\,(T-T_\infty)\quad\text{am Rand.}
\end{gathered}
\label{eq:fourier}
\end{equation}
\emph{In Worten:} Ein Volumenelement erwärmt sich (links), wenn Wärme von
Nachbarn zufließt oder vor Ort entsteht (rechts). Am Rand gibt es umso mehr
Wärme ab, je wärmer es als die Luft ist.
\py{\url{thermal_solver.py} (Anhang~B.3)}

\subsubsection{Temperaturabhängigkeit}
Warme Metalle leiten schlechter:
\begin{equation}
\sigma(T) = \frac{\sigma_0}{1+\alpha (T-T_0)},\qquad
\alpha \approx 3{,}9\cdot10^{-3}\ \mathrm{K^{-1}}.
\label{eq:sigma}
\end{equation}
Die Wirkung ist gegenläufig: Die Wirbelstromverluste der Scheibe sinken, die
Spulenverluste steigen.

\begin{table}[t]
\centering\footnotesize
\caption{Das Modell auf einen Blick}
\label{tab:model}
\begin{tabularx}{\columnwidth}{@{}lYYY@{}}
\toprule
Gl. & liefert & braucht & Kopplung \\
\midrule
\eqref{eq:em} Magnetfeld & $A_\varphi(r,z)$ & Geometrie, Material, Strom & $\rightarrow$ \eqref{eq:loss} \\
\eqref{eq:loss} Wärmequelle & $q(r,z)$, Leistung je Bauteil & $A_\varphi$ & $\rightarrow$ \eqref{eq:fourier} \\
\eqref{eq:fourier} Wärmeleitung & $T(r,z,t)$ & $q$, Wärmedaten, $h$ & $\rightarrow$ \eqref{eq:sigma} \\
\eqref{eq:sigma} Leitfähigkeit & Korrekturfaktor & Temperatur & $\rightarrow$ \eqref{eq:em}, langsam \\
\bottomrule
\end{tabularx}
\end{table}


%%%%%%%%%%%%%%%%%%%%%%%%%%%%%%%%%%%%%%%%%%%%%%%%%%%%%%%%%%%%%%%%%%%%%%%%%%%%%%%
\section{Methodik}
\label{sec:method}

\subsection{Überblick}

Die Gleichungen aus Abschnitt~\ref{sec:problem} werden in \textbf{zwei
Stufen} gelöst (Abbildung~\ref{fig:stages}):
\begin{itemize}[nosep]
  \item \textbf{Offline (einmal):} Die FEM löst
        Gl.~\eqref{eq:em}--\eqref{eq:fourier} genau. Das dauert etwa eine
        Sekunde und braucht Löserbibliotheken.
  \item \textbf{Online (laufend):} Das daraus abgeleitete reduzierte Modell (ROM)
        rechnet jeden Zeitschritt in Mikrosekunden, nur mit Grundrechenarten.
\end{itemize}

\begin{figure}[t]
\centering
\begin{tikzpicture}[font=\scriptsize, node distance=2.5mm,
    box/.style={draw, rounded corners, align=center, text width=12.5mm,
                minimum height=11mm, inner sep=1.5pt},
    arr/.style={-Stealth}]
  \node[box] (A) {Gleichungen \eqref{eq:em}--\eqref{eq:sigma}};
  \node[box, right=of A, fill=blue!8] (B) {FEM\\genau, langsam\\einmal};
  \node[box, right=of B] (C) {Reduktion\\A, B, C};
  \node[box, right=of C, fill=green!10] (D) {ROM\\8 Zustände\\jeder Schritt};
  \node[box, right=of D] (E) {Temperatur, Schwebe\-höhe\\HTML};
  \draw[arr] (A) -- (B); \draw[arr] (B) -- (C); \draw[arr] (C) -- (D); \draw[arr] (D) -- (E);
\end{tikzpicture}
\caption{Zwei Stufen: offline einmal genau rechnen, online klein und schnell.}
\label{fig:stages}
\end{figure}

Abschnitt~\ref{sec:fem} erklärt die FEM, Abschnitt~\ref{sec:mor} die
Reduktion.

\subsection{Finite-Elemente-Methode}
\label{sec:fem}

\subsubsection{Numerisches Verfahren}
Für ein Gerät mit fünf Materialbereichen gibt es keine Formellösung. Man
zerlegt deshalb das Gebiet in kleine Stücke und löst die Gleichung
stückweise näherungsweise. Die FEM ist dafür das Standardwerkzeug der
TEAM-Benchmarks~\cite{team28,biro1989vectorpotential}: Jedes Stück trägt
seinen eigenen Materialwert, und Feld- und Wärmeproblem lassen sich mit
demselben Programmgerüst lösen. Ein Vergleich mit anderen Verfahren steht in
Anhang~C.1.

\subsubsection{Diskretisierung}
Die Ebene $(r,z)$ wird in kleine \textbf{Dreiecke} zerlegt (Netz). In jedem
Dreieck ändert sich die gesuchte Größe linear. Sie ist dann durch ihre Werte
an den Ecken (\textbf{Knoten}) festgelegt. Aus einer unbekannten Funktion
werden so endlich viele unbekannte Zahlen.

Die Gleichung wird nicht in jedem Punkt, sondern \textbf{im Mittel um jeden
Knoten} erfüllt. Physikalisch ist das eine Wärmebilanz um jeden Knoten
(\emph{schwache Form}):
\begin{multline}
\int_\Omega k\nabla T\cdot\nabla v_i\,2\pi r\,dA + \oint_\Gamma hTv_i\,2\pi r\,ds \\
= \int_\Omega q\,v_i\,2\pi r\,dA+\oint_\Gamma hT_\infty v_i\,2\pi r\,ds .
\label{eq:weak}
\end{multline}
\emph{In Worten:} Wärmeleitung + Abgabe am Rand = erzeugte Wärme. Der Faktor
$2\pi r$ berücksichtigt die Drehung um die Achse. $v_i$ ist eine
"`Zeltfunktion"', die am Knoten $i$ gleich 1 ist und zu den Nachbarn auf 0
abfällt.
\py{Elementformeln in Anhang~D}

\subsubsection{Gleichungssystem}
Schreibt man Gl.~\eqref{eq:weak} für alle Knoten auf, entsteht ein lineares
Gleichungssystem:
\begin{equation}
\mathbf K\,\mathbf T=\mathbf f,\qquad\text{zeitabhängig:}\quad
\mathbf M\,\dot{\mathbf T}+\mathbf K\,\mathbf T=\mathbf f .
\label{eq:lin}
\end{equation}
$\mathbf T$ enthält alle Knotentemperaturen, $\mathbf K$ die Leitung,
$\mathbf M$ die Wärmekapazität, $\mathbf f$ die Quellen. Fast alle Einträge
von $\mathbf K$ sind null, weil nur Nachbarknoten gekoppelt sind. Für
Gl.~\eqref{eq:em} läuft alles gleich, nur mit komplexen Zahlen.
\py{\url{scipy.sparse.linalg.spsolve}~\cite{virtanen2020scipy} (Anhang~B.2/B.3)}

\subsubsection{Modellgröße}
Zwei physikalische Vereinfachungen halten das Problem klein:
\begin{itemize}[nosep]
  \item \textbf{Phasor statt Zeitschritt:} Das Feld schwingt 50-mal pro
        Sekunde, die Temperatur ändert sich über Minuten. Die Wärme spürt nur
        den Mittelwert, und genau den liefert Gl.~\eqref{eq:loss}.
  \item \textbf{Grobes Netz in Dickenrichtung:} Das Feld dringt in Aluminium
        bei 50\,Hz etwa 12\,mm tief ein, die Scheibe ist nur 3\,mm dick. Der
        Wirbelstrom ist über die Dicke fast gleich.
\end{itemize}
Damit hat das Feldmodell rund 8700, das Wärmemodell der Scheibe rund 700
Unbekannte.

\subsubsection{Plausibilitätsprüfung}
Eine FEM liefert immer Zahlen, auch bei Fehlern. Deshalb prüft der Code bei
jedem Lauf Größen, deren richtiges Ergebnis bekannt ist:
\begin{enumerate}[nosep]
  \item \textbf{Energiebilanz:} Erzeugte Wärme = abgegebene Wärme. Ergebnis:
        Übereinstimmung auf Rundungsgenauigkeit.
  \item \textbf{Gebietsgröße:} Zwei verschiedene Randbedingungen am
        Außenrand ergeben nahezu dieselben Verluste (Abweichung $<1\,\%$). Das
        Rechengebiet ist also groß genug.
  \item \textbf{Linearität:} Doppelter Strom ergibt vierfache Verluste.
\end{enumerate}

\subsection{Modellordnungsreduktion (MOR)}
\label{sec:mor}

\subsubsection{Grundidee}
Die FEM ist genau, aber für jeden neuen Strom und jeden Zeitschritt müsste
sie neu gelöst werden. Die Reduktion nutzt aus, dass die Tausenden
Unbekannten \textbf{nicht unabhängig} sind: Beim Aufheizen ändert sich vor
allem, \emph{wie warm} die Scheibe ist, kaum aber, \emph{wo} sie wärmer ist.

\begin{remark}[Analogie]
Ein Foto, das nur heller oder dunkler wird. Man speichert das Bild einmal und
merkt sich pro Moment nur einen Helligkeitsfaktor.
\end{remark}

Mathematisch wird das Feld als Summe weniger fester "`Bilder"'
$\mathbf v_k$ mit zeitabhängigen Gewichten $\beta_k(t)$
geschrieben~\cite{benner2015survey}:
\begin{equation}
\mathbf T(t)\approx T_\infty+\sum_{k=1}^{r}\beta_k(t)\,\mathbf v_k,\qquad r \ll N .
\label{eq:rb}
\end{equation}
Drei Reduktionen setzen diese Idee um (Tabelle~\ref{tab:mor}). Sie sind
keine Alternativen, sondern wirken gleichzeitig, jede auf einen anderen Teil
des Geräts (anschaulich: Anhang~F).

\subsubsection{Reduktion A: Strom-Skalierung}
Gl.~\eqref{eq:em} ist \textbf{linear}: Doppelter Strom ergibt doppeltes
Feld, an jedem Ort gleich. Nach Gl.~\eqref{eq:loss} vervierfacht sich dann
die Wärmequelle, ebenfalls überall gleich. Die \emph{Form} der
Verlustverteilung bleibt, nur ihre Höhe ändert sich. Die Feld-FEM wird daher
nur \textbf{einmal} bei $I_{\mathrm{ref}}=5\,\mathrm A$ gelöst. Für jeden
Strom gilt dann:
\begin{equation}
q(r,z;I,T)=\hat q(r,z)\cdot\Big(\frac{I}{I_{\mathrm{ref}}}\Big)^{2}\cdot
\frac{\sigma(T)}{\sigma(T_{\mathrm{ref}})} .
\label{eq:realtime}
\end{equation}
Gültig ist das, solange das Eisen nicht magnetisch sättigt. Am Betriebspunkt
ist das erfüllt; bei höheren Strömen wird es knapp (Abschätzung in
Anhang~C.4).
\py{\url{rom.py} und \url{twin_core.py} (Anhang~B.4/B.5)}

\subsubsection{Reduktion B: eine Temperaturform für die Scheibe}
Aluminium leitet Wärme sehr gut, die Scheibe ist dünn und gibt Wärme nur
langsam an die Luft ab. Die Scheibe ist daher fast überall gleich warm. Als
einziges "`Bild"' dient das stationäre FEM-Temperaturfeld
$\Delta T_{\mathrm{ref}}(r,z)$. Aus Gl.~\eqref{eq:rb} wird eine einzige
Gleichung für das Gewicht $\beta(t)$:
\begin{equation}
\tau\,\dot\beta=\Big(\frac{I}{I_{\mathrm{ref}}}\Big)^{2}s-\beta,\qquad
T(r,z,t)=T_\infty+\beta(t)\,\Delta T_{\mathrm{ref}}(r,z).
\label{eq:beta}
\end{equation}
\emph{In Worten:} $\beta=0$ heißt kalt, $\beta=1$ so warm wie im
Referenzfall. $\beta$ strebt dem Zielwert $(I/I_{\mathrm{ref}})^2 s$ zu ($s$:
Leitfähigkeitsfaktor aus Gl.~\eqref{eq:sigma}). Die Zeitkonstante $\tau$
(rund 4 Minuten) gibt das Tempo an. Dass eine einzige Form genügt, ist durch
eine Eigenwertanalyse und einen Vergleich mit der vollen FEM belegt
(Anhang~C.2, C.3).
\py{\url{rom.py} (Anhang~B.4)}

\subsubsection{Reduktion C: RC-Netzwerk für Spulen und Eisen}
Spulen und Eisen werden durch wenige Temperaturknoten beschrieben, wie ein
elektrischer Schaltkreis: Wärmekapazität $\widehat{=}$ Kondensator,
Wärmeleitwert $\widehat{=}$ Widerstand~\cite{wallscheid2021thermal}. Für
jeden Knoten $i$ gilt:
\begin{equation}
C_i\,\dot T_i=P_i(I)+\sum_j G_{ij}(T_j-T_i)-hA_i\,(T_i-T_{\mathrm{Luft}}).
\label{eq:rc}
\end{equation}
\emph{In Worten:} Ein Knoten wird wärmer durch Verlustleistung und Wärme von
wärmeren Nachbarn, er kühlt durch Abgabe an die Luft.

Das Netz hat sechs Knoten (Abbildung~\ref{fig:rc}). Jede Spule hat zwei: die
sichtbare Oberfläche und das träge Wicklungsinnere. Dazu kommen Eisen und
die Luft im Gerät. Die Zahlenwerte stehen in Anhang~A.
\py{\url{twin_core.py} (Anhang~B.5)}

\begin{figure}[t]
\centering
\begin{tikzpicture}[font=\scriptsize,
    nd/.style={circle, draw, fill=white, minimum size=8.5mm, inner sep=0pt},
    src/.style={draw, rounded corners=5pt, fill=orange!20, inner sep=2pt},
    amb/.style={draw, fill=gray!15, inner sep=3pt},
    arr/.style={-Stealth}, darr/.style={Stealth-Stealth}]
  \node[nd] (ID) at (0,2.2)   {iW};
  \node[nd] (IS) at (1.6,2.2) {iO};
  \node[nd, fill=gray!25] (FE) at (3.3,2.2) {Fe};
  \node[nd] (OS) at (5.0,2.2) {aO};
  \node[nd] (OD) at (6.6,2.2) {aW};
  \node[src] (Pi) at (1.6,3.5) {$P_{\text{innen}}$};
  \node[src] (Pf) at (3.3,3.5) {$P_{\text{Eisen}}$};
  \node[src] (Po) at (5.0,3.5) {$P_{\text{außen}}$};
  \node[nd, fill=cyan!10] (AIR) at (3.3,0.9) {L};
  \node[amb] (AMB) at (3.3,-0.35) {Umgebung};
  \draw[arr] (Pi) -- (IS); \draw[arr] (Pf) -- (FE); \draw[arr] (Po) -- (OS);
  \draw[darr] (ID) -- (IS); \draw[darr] (OD) -- (OS); \draw[darr] (IS) -- (FE);
  \draw[arr] (IS) -- (AIR); \draw[arr] (OS) -- (AIR); \draw[arr] (FE) -- (AIR);
  \draw[arr] (AIR) -- (AMB);
\end{tikzpicture}
\caption{Thermisches RC-Netzwerk. iO/iW: Innenspule Oberfläche/Wicklung,
aO/aW: Außenspule Oberfläche/Wicklung, Fe: Eisen, L: Luft im Gerät,
$P$: Verluste. Pfeile: Wärmefluss.}
\label{fig:rc}
\end{figure}

\subsubsection{Kalibrierung}
\label{sec:calib}
Die zwei wichtigsten Parameter, die Wärmeübergänge der Spulen
$hA_{\text{innen}}$ und $hA_{\text{außen}}$, werden so eingestellt, dass
das Modell die IR-Messung bei $7{,}8\,\mathrm A$ trifft
($79$ / $74\,^\circ$C). Ein Suchverfahren variiert beide Werte, bis das
\textbf{laufende} Modell genau diese Temperaturen erreicht.
\py{\url{refit_hA.py} (Anhang~B.6)}

\subsubsection{Ergebnis: das ROM}
Aus mehreren Tausend Unbekannten werden \textbf{acht Zustandsgrößen}
(Tabelle~\ref{tab:mor}); das ist das ROM, im Code \url{rom.py} und
\url{twin_core.py}. Ein Zeitschritt ist reine Arithmetik und dauert
Mikrosekunden statt Sekunden.

\begin{table*}[t]
\centering\footnotesize
\caption{Die drei Reduktionen}
\label{tab:mor}
\begin{tabularx}{\textwidth}{@{}lYYYY@{}}
\toprule
 & Grund & vorher & nachher & Grenze \\
\midrule
A: Strom-Skalierung & Gl.~\eqref{eq:em} ist linear & 8700 komplexe Unbekannte & 1 Multiplikation & Sättigung des Eisens \\
B: eine Form & Scheibe fast gleich warm & ca.\ 700 Unbekannte & 2 Gewichte $\beta$ & Anhang~C.3 \\
C: RC-Netz & Spulen/Eisen kompakt & -- & 6 Knoten & Parameter teils geschätzt \\
\bottomrule
\end{tabularx}
\end{table*}


%%%%%%%%%%%%%%%%%%%%%%%%%%%%%%%%%%%%%%%%%%%%%%%%%%%%%%%%%%%%%%%%%%%%%%%%%%%%%%%
\section{Berechnungswerkzeug}
\label{sec:code}

Der Code ist in Python geschrieben (NumPy, SciPy~\cite{virtanen2020scipy}).
Alle Zahlenwerte stehen in einer einzigen Datei, \url{params.yaml}.
Abbildung~\ref{fig:pipeline} zeigt den Datenfluss. Was jedes Modul
einliest, wie es rechnet und welches Verfahren es nutzt, erklärt Anhang~B in
einfachen Worten.

\begin{figure}[t]
\centering
\begin{tikzpicture}[font=\scriptsize, node distance=3.2mm,
    mod/.style={draw, rounded corners, fill=white, align=center,
                text width=31mm, inner sep=2pt},
    arr/.style={-Stealth}]
  \node[mod, fill=yellow!15] (P) {\texttt{params.yaml}\\Geometrie, Material, Messwerte};
  \node[mod, below=of P] (C) {\texttt{config.py}\\Werte laden};
  \node[mod, below=of C] (E) {\texttt{em\_solver.py}\\Feld-FEM, Gl.~\eqref{eq:em}--\eqref{eq:loss}};
  \node[mod, below=of E] (T) {\texttt{thermal\_solver.py}\\Wärme-FEM, Gl.~\eqref{eq:fourier}};
  \node[mod, below=of T] (R) {\texttt{rom.py}\\Reduktion A + B};
  \node[mod, right=5mm of E] (L) {\texttt{build\_twin\_html\_}\\\texttt{fem.py}\\Reduktion C};
  \node[mod, right=5mm of C] (K) {\texttt{refit\_hA.py}\\Kalibrierung $hA$};
  \node[mod, below=7mm of R, fill=green!10] (TC) {\texttt{twin\_core.py}\\Gl.~\eqref{eq:beta} + \eqref{eq:rc}};
  \node[mod, below=of TC, fill=cyan!10] (V) {\texttt{digital\_twin\_fem.html}\\Ansicht im Browser};
  \begin{scope}[on background layer]
    \node[draw, dashed, fill=blue!4, inner sep=2.5mm, fit=(C)(E)(T)(R)(L)(K),
          label={[anchor=south east, font=\scriptsize\bfseries]south east:OFFLINE -- einmal}] {};
    \node[draw, dashed, fill=green!4, inner sep=2mm, fit=(TC),
          label={[anchor=north west, font=\scriptsize\bfseries]south east:ONLINE -- jeder Zeitschritt}] {};
  \end{scope}
  \draw[arr] (P) -- (C); \draw[arr] (C) -- (E); \draw[arr] (E) -- (T);
  \draw[arr] (T) -- (R); \draw[arr] (R) -- (TC); \draw[arr] (TC) -- (V);
  \draw[arr] (E) -- (L); \draw[arr] (L) |- (TC);
  \draw[arr, dashed] (K) |- node[pos=0.25, right, font=\tiny] {schreibt $hA$} (P);
\end{tikzpicture}
\caption{Datenfluss durch den Code.}
\label{fig:pipeline}
\end{figure}

\textbf{Starten.} \url{RUN.py} im Projektordner (bzw.\ Doppelklick auf
\url{RUN.command}) rechnet bei Bedarf neu und öffnet die HTML-Ansicht im
Browser.

\textbf{Stromeingang.} Standardmäßig ist der Strom fest ($5\,\mathrm A$)
oder über einen Schieberegler einstellbar. Ein Stromsensor (ACS712 am
Arduino) ist softwareseitig vorbereitet (\url{data_io.py}); der Einbau am
Versuchsstand ist in Arbeit.


%%%%%%%%%%%%%%%%%%%%%%%%%%%%%%%%%%%%%%%%%%%%%%%%%%%%%%%%%%%%%%%%%%%%%%%%%%%%%%%
\section{Ergebnisse}
\label{sec:results}

\subsection{Verlustleistungen}

\begin{table}[t]
\centering\footnotesize
\caption{Verluste aus der Feld-FEM ($5\,\mathrm A$, $20\,^\circ$C)}
\label{tab:losses}
\begin{tabular}{@{}lrr@{}}
\toprule
Quelle & $P$ [W] & Anteil \\
\midrule
Wirbelstrom Scheibe & 25{,}81 & 19\,\% \\
Wirbelstrom Eisen & 5{,}86 & 4\,\% \\
Innenspule & 52{,}32 & 38\,\% \\
Außenspule & 54{,}12 & 39\,\% \\
\midrule
\textbf{Gesamt} & \textbf{138{,}11} & 100\,\% \\
\bottomrule
\end{tabular}
\end{table}

Die Spulen erzeugen rund drei Viertel der Wärme
(Tabelle~\ref{tab:losses}). Das passt dazu, dass sie im Wärmebild die
heißesten Bauteile sind. (Zur Stromkonvention Effektivwert/Amplitude:
Anhang~C.5.)

\subsection{Vergleich mit der Messung}

\begin{table*}[t]
\centering\footnotesize
\caption{Messwert gegen Modell}
\label{tab:compare}
\begin{tabular}{@{}lrrrl@{}}
\toprule
Messgröße & Messung & Modell & Differenz & Bewertung \\
\midrule
Innenspule, stationär, 7{,}8\,A & 79\,$^\circ$C & 79{,}0\,$^\circ$C & 0 & kalibriert, kein Nachweis \\
Außenspule, stationär, 7{,}8\,A & 74\,$^\circ$C & 74{,}0\,$^\circ$C & 0 & kalibriert, kein Nachweis \\
Außenspule, Rampe, $t = 140$\,s & 40{,}0\,$^\circ$C & 33{,}6\,$^\circ$C & $-6{,}5$\,K & Vorhersage \\
Außenspule, Rampe, $t = 200$\,s & 43{,}0\,$^\circ$C & 34{,}7\,$^\circ$C & $-8{,}3$\,K & Vorhersage \\
Außenspule, Rampe, $t = 300$\,s & 48{,}5\,$^\circ$C & 36{,}2\,$^\circ$C & $-12{,}3$\,K & Vorhersage \\
Außenspule, Rampe, $t = 360$\,s & 55{,}8\,$^\circ$C & 40{,}0\,$^\circ$C & $-15{,}8$\,K & Vorhersage \\
Außenspule, Rampe, $t = 450$\,s & 56{,}0\,$^\circ$C & 43{,}6\,$^\circ$C & $-12{,}4$\,K & Vorhersage \\
Scheibe, stationär, 7{,}8\,A & 37--40\,$^\circ$C & 107\,$^\circ$C & $\approx +68$\,K & IR unzuverlässig \\
Schwebehöhe (sichtbar), 5\,A & 7--8\,mm & 14{,}7\,mm & $\approx +7$\,mm & Vorhersage \\
\bottomrule
\end{tabular}
\end{table*}

\textbf{Kurz zusammengefasst} (Tabelle~\ref{tab:compare}):
\begin{itemize}[nosep]
  \item \textbf{Spulen, stationär:} exakt, aber nur, weil darauf kalibriert
        wurde.
  \item \textbf{Spulen, Zeitverlauf:} Das Modell heizt zu langsam auf
        (mittlerer Fehler $11{,}5\,\mathrm K$).
  \item \textbf{Scheibe:} Die IR-Kamera misst auf blankem Aluminium zu tief;
        ein Urteil ist nicht möglich.
  \item \textbf{Schwebehöhe:} Das Modell überschätzt den Spalt um etwa
        $7\,\mathrm{mm}$.
\end{itemize}

\subsection{Ursachen der Abweichungen}
\label{sec:causes}

\begin{itemize}[nosep]
  \item \textbf{Zeitverlauf zu langsam:} Die Endtemperatur hängt nur von den
        kalibrierten Wärmeübergängen ab. Das Tempo hängt dagegen von den
        Wärmekapazitäten und der Aufteilung Oberfläche/Wicklung ab. Diese
        lassen sich aus einer Aufheizkurve allein nicht bestimmen; dafür
        fehlt eine \textbf{Abkühlkurve}.
  \item \textbf{Scheibe:} Blankes Aluminium strahlt wenig Wärme ab und
        spiegelt die Umgebung. Die Kamera zeigt daher zu tiefe Werte.
        Abhilfe: ein mattschwarzer Messfleck.
  \item \textbf{Schwebehöhe:} Wahrscheinlichste Ursache ist die angenommene
        Eisenpermeabilität $\mu_r = 1000$. Sie wurde bewusst \emph{nicht}
        nachträglich an die Beobachtung angepasst. Zudem rechnet das
        Feldmodell mit fester Scheibenhöhe.
\end{itemize}


%%%%%%%%%%%%%%%%%%%%%%%%%%%%%%%%%%%%%%%%%%%%%%%%%%%%%%%%%%%%%%%%%%%%%%%%%%%%%%%
\section{Ausblick}
\label{sec:outlook}

Fast alle Abweichungen haben dieselbe Ursache: Es fehlen Messdaten. Die
nächsten Schritte sind daher:
\begin{enumerate}[nosep]
  \item \textbf{Abkühlkurve} der Spulen aufnehmen $\rightarrow$ Zeitverlauf
        kalibrieren;
  \item \textbf{mattschwarzer Messfleck} auf Scheibe und Eisen
        $\rightarrow$ verlässliche IR-Werte;
  \item \textbf{Eisenpermeabilität} messen $\rightarrow$ Schwebehöhe klären;
  \item \textbf{Stromsensor} einbauen $\rightarrow$ Modell läuft mit dem
        gemessenen Strom statt mit einem festen Wert;
  \item \textbf{unabhängiger Messpunkt bei 5\,A} $\rightarrow$ Kalibrierung
        prüfen, nicht nur anpassen.
\end{enumerate}
Das Vorgehen -- FEM einmal lösen, Reduzierbarkeit prüfen, online das ROM
rechnen -- lässt sich auf 3D-Modelle
übertragen~\cite{hartmann2025golden,hartmann2018mor}.


%%%%%%%%%%%%%%%%%%%%%%%%%%%%%%%%%%%%%%%%%%%%%%%%%%%%%%%%%%%%%%%%%%%%%%%%%%%%%%%%
% Anhänge mit Buchstaben (A, B, ...) statt der Standard-Römischen Zahlen, damit
% die Punkte A.1, B.2, ... eindeutig sind (Schalter aus ieeeconf.cls).
\useRomanappendicesfalse
\appendices

\section{Mess- und Rechenwerte}
\label{app:values}

% Überschrift + Tabelle als ein Block (minipage), damit sie nicht getrennt werden.
\noindent\begin{minipage}{\columnwidth}
\anh{A.1 Messwerte und Annahmen (Eingang ins Modell)}

\begin{center}\footnotesize
\begin{tabularx}{\columnwidth}{@{}>{\raggedright\arraybackslash}p{2.1cm}Y>{\raggedright\arraybackslash}p{2.0cm}@{}}
\toprule
Größe & Wert & Herkunft \\
\midrule
Strom Betrieb / Kalibrierung & 5{,}0 / 7{,}8\,A eff., 50\,Hz & gemessen (Multimeter) \\
Variac-Stellung & $220^\circ \approx 195$\,V $\rightarrow$ 5\,A; $270^\circ$ (max.) $\approx 240$\,V $\rightarrow$ 7{,}8\,A & gemessen (Stellwinkel) \\
Radien Kern / Innenspule / Ring / Außenspule & 0--25{,}9 / 27{,}9--61{,}9 / 64{,}9--79{,}9 / 82{,}9--102{,}9\,mm & gemessen (Lineal, 10.07.2026) \\
Höhe Kern/Spulen/Ring & 53\,mm & gemessen \\
Scheibe $R$ / $d$ / Masse & 80\,mm / 3\,mm / 159--163\,g (Kraftrechnung: 162{,}9\,g) & gemessen \\
Windungszahl innen / außen & ca.\ 1000 / 500 & Angabe zum Aufbau \\
$\sigma$ Al / Cu & 34 / 59{,}6\,MS/m & Literatur \\
Al: $k$, $\rho$, $c_p$, $\alpha$ & 237\,W/(m\,K), 2700\,kg/m$^3$, 900\,J/(kg\,K), $3{,}9\cdot10^{-3}$\,1/K & Literatur \\
Eisen $\mu_r$ / $\sigma$ / $B_{\mathrm{sat}}$ & 1000 / 1\,MS/m / 1{,}5\,T & angenommen \\
Draht / Füllfaktor & 1{,}2\,mm / 0{,}6 & angenommen \\
$h$ Scheibe oben / unten & 10 / 25\,W/(m$^2$\,K) & angenommen \\
Luftknoten $C$ / $hA$ & 3000\,J/K / 40\,W/K & angenommen \\
$hA$ innen / außen & 2{,}474 / 2{,}889\,W/K & kalibriert \\
IR stationär 7{,}8\,A, $29\,^\circ$C & Innenspule 79, Außenspule 74, Kern 45, Scheibe 37--40\,$^\circ$C & gemessen (IR, 23.06.2026) \\
IR Rampe Außenspule & 40 / 43 / 48{,}5 / 55{,}75 / 56\,$^\circ$C bei 140 / 200 / 300 / 360 / 450\,s & gemessen (Zeit $\pm$10\,s) \\
Schwebehöhe sichtbar, 5\,A & 7--8\,mm & beobachtet \\
\bottomrule
\end{tabularx}
\end{center}
\end{minipage}

\noindent\begin{minipage}{\columnwidth}
\anh{A.2 Rechenwerte (Ausgabe des Modells)} Kennung = Zeile in
\url{report_numbers_output.txt}.

\begin{center}\footnotesize
\begin{tabularx}{\columnwidth}{@{}lYY@{}}
\toprule
Kenn. & Größe & Wert \\
\midrule
E1 & Feld-FEM: Knoten / Dreiecke / Rechenzeit & 8700 / 17\,028 / 0{,}40\,s \\
E2 & Verlustrechnung inkl.\ Eisenkorrektur & 1{,}3\,s \\
E3, E4 & Verluste Scheibe / Eisen / Innen / Außen & 25{,}81 / 5{,}86 / 52{,}32 / 54{,}12\,W \\
E6 & Verlustverhältnis bei doppeltem Strom & 4{,}000 (linear), 4{,}009 mit Eisenkorrektur \\
E7 & Gebietsgröße, größte Abweichung & 0{,}77\,\% \\
T1 & Wärme-FEM: Knoten / Dreiecke / Rechenzeit & 697 / 1280 / 0{,}03\,s \\
T2 & Energiebilanz & 25{,}81388\,W = 25{,}81388\,W \\
R1 & Scheibe $C$ / $UA$ / $\tau$ & 146{,}6\,J/K / 0{,}600\,W/K / 244\,s \\
L1 & Wärmekapazitäten Spule Oberfläche / Wicklung & innen 247 / 854, außen 255 / 884\,J/K; Eisen 1676\,J/K \\
L3 & stationär 7{,}8\,A / $29\,^\circ$C & Spule 79{,}0 / 74{,}0; Eisen 83{,}7; Wicklung 178 / 176; Scheibe 107\,$^\circ$C \\
L3 & stationär 5\,A / $20\,^\circ$C & Spule 44{,}0 / 41{,}5; Eisen 46{,}3; Wicklung 85 / 84; Scheibe 58\,$^\circ$C \\
L4, L5 & Rampe & Tabelle~\ref{tab:compare}; RMS 11{,}5\,K \\
L6 & ein Online-Zeitschritt & $\approx 3\,\mu$s \\
V1 & Schwebehöhe & 11{,}7\,mm (Unterkante), 14{,}7\,mm sichtbar \\
\bottomrule
\end{tabularx}
\end{center}
\end{minipage}

\smallskip\noindent Hinweis: Die Wicklungstemperatur von $178\,^\circ$C bei
$7{,}8\,\mathrm A$ ist nicht gemessen und vermutlich zu hoch, weil der
innere Wärmeleitwert nur geschätzt ist (vgl.\
Abschnitt~\ref{sec:causes}).

%%%%%%%%%%%%%%%%%%%%%%%%%%%%%%%%%%%%%%%%%%%%%%%%%%%%%%%%%%%%%%%%%%%%%%%%%%%%%%%%
\section{Umsetzung im Code}
\label{app:code}

\begingroup\raggedright
Jedes Modul wird nach demselben Schema beschrieben: \textbf{Was kommt rein?
Was passiert? Mit welchem Verfahren?}

\anh{B.1 \texttt{config.py} -- Werte laden}
\begin{itemize}[nosep]
  \item \emph{Rein:} die Datei \url{params.yaml}.
  \item \emph{Was passiert:} Alle Zahlen werden gelesen, Millimeter in Meter
        umgerechnet und abgeleitete Werte gebildet, z.\,B.\ Spitzenstrom =
        Effektivwert $\cdot\sqrt2$.
  \item \emph{Verfahren:} einfaches Einlesen (Bibliothek PyYAML); keine
        Rechnung.
\end{itemize}

\anh{B.2 \texttt{em\_solver.py} -- Magnetfeld und Wärmequelle (Gl.~\eqref{eq:em}, \eqref{eq:loss})}
\begin{itemize}[nosep]
  \item \emph{Rein:} Geometrie, Materialwerte, Strom.
  \item \emph{Was passiert:} (1) Das Gebiet ($1\,\mathrm m \times
        1\,\mathrm m$) wird in Dreiecke zerlegt, fein am Gerät, grob weiter
        weg. (2) Jedes Dreieck bekommt seine Materialwerte. (3) Für jedes
        Dreieck wird ein kleiner Beitrag zur großen Matrix berechnet und
        eingetragen. (4) Das Gleichungssystem wird gelöst. (5) Das Eisen wird
        nachkorrigiert, falls das Feld stark ist. (6) Aus dem Feld wird
        Gl.~\eqref{eq:loss} berechnet und über jedes Bauteil aufsummiert.
  \item \emph{Verfahren:} FEM mit linearen Dreiecken; Lösen mit
        \url{scipy.sparse.linalg.spsolve} (direktes Verfahren, ähnlich dem
        Gauß-Verfahren aus der Schule, aber für sehr große, fast leere
        Matrizen).
  \item \emph{Raus:} Verlustkarte der Scheibe, Leistung je Bauteil,
        Hubkraft.
\end{itemize}

\anh{B.3 \texttt{thermal\_solver.py} -- Temperatur der Scheibe (Gl.~\eqref{eq:fourier}, \eqref{eq:weak})}
\begin{itemize}[nosep]
  \item \emph{Rein:} Verlustkarte aus B.2, Wärmedaten, Wärmeübergänge.
  \item \emph{Was passiert:} Die Verlustkarte wird auf ein feineres Netz der
        Scheibe übertragen, Gl.~\eqref{eq:weak} als Gleichungssystem
        aufgestellt und gelöst. Danach wird geprüft, ob erzeugte =
        abgegebene Wärme.
  \item \emph{Verfahren:} FEM, \url{spsolve}, lineare Interpolation.
  \item \emph{Raus:} stationäres Temperaturfeld der Scheibe.
\end{itemize}

\anh{B.4 \texttt{rom.py} -- Reduktion A und B (Gl.~\eqref{eq:realtime}, \eqref{eq:beta})}
\begin{itemize}[nosep]
  \item \emph{Rein:} Temperaturfeld aus B.3.
  \item \emph{Was passiert:} Das Temperaturfeld wird als "`Bild"'
        $\Delta T_{\mathrm{ref}}$ gespeichert. Aus Masse und Wärmeabgabe
        werden Wärmekapazität $C$, Abgabe $UA$ und Zeitkonstante
        $\tau = C/UA$ berechnet.
  \item \emph{Verfahren:} Energiebilanz eines einzelnen Körpers; zur
        Kontrolle numerische Zeitintegration
        (\url{scipy.integrate.solve_ivp}).
  \item \emph{Raus:} $\Delta T_{\mathrm{ref}}$, $\tau$.
\end{itemize}

\anh{B.5 \texttt{twin\_core.py} -- das Online-Modell (Gl.~\eqref{eq:beta}, \eqref{eq:rc})}
\begin{itemize}[nosep]
  \item \emph{Rein:} aktueller Strom $I$, Zeitschritt $\Delta t$, alle
        Koeffizienten aus B.4 und dem RC-Netz.
  \item \emph{Was passiert pro Zeitschritt:} (1) Verluste mit
        $(I/I_{\mathrm{ref}})^2$ skalieren. (2) Jede Temperatur um
        "`Änderungsrate $\times$ Zeitschritt"' weiterschieben. (3) Ist der
        Zeitschritt zu groß, wird er in kleinere Teilschritte zerlegt, damit
        die Rechnung stabil bleibt. (4) Schwebehöhe mit einer fertigen Formel
        für eine gedämpfte Schwingung berechnen.
  \item \emph{Verfahren:} explizites Euler-Verfahren ("`neuer Wert = alter
        Wert + Steigung $\times$ Schrittweite"'); nur NumPy und
        Standardbibliothek.
  \item \emph{Raus:} Temperaturen aller Knoten, Temperaturfeld der Scheibe,
        Schwebehöhe.
\end{itemize}

\anh{B.6 \texttt{refit\_hA.py} -- Kalibrierung}
\begin{itemize}[nosep]
  \item \emph{Rein:} gemessene Spulentemperaturen ($79$ / $74\,^\circ$C bei
        $7{,}8\,\mathrm A$).
  \item \emph{Was passiert:} Das Programm probiert Werte für
        $hA_{\text{innen}}$ und $hA_{\text{außen}}$, lässt das
        Online-Modell jeweils bis zum Gleichgewicht laufen und vergleicht
        mit der Messung, bis die Differenz null ist.
  \item \emph{Verfahren:} Nullstellensuche (\url{scipy.optimize.fsolve}).
  \item \emph{Raus:} kalibrierte $hA$-Werte in \url{params.yaml}.
\end{itemize}

\anh{B.7 \texttt{build\_twin\_html\_fem.py} -- Reduktion C und HTML-Ansicht}
\begin{itemize}[nosep]
  \item \emph{Rein:} Ergebnisse aus B.2--B.4, Geometrie.
  \item \emph{Was passiert:} Die Parameter des RC-Netzes werden berechnet
        (Kapazitäten aus Volumen und Werkstoff). Alle Koeffizienten und das
        Online-Modell werden als JavaScript in eine HTML-Datei geschrieben,
        damit die Ansicht ohne Python im Browser läuft.
  \item \emph{Raus:} \url{outputs/digital_twin_fem.html}.
\end{itemize}

\anh{B.8 \texttt{data\_io.py} -- Stromeingang (vorbereitet)}
\begin{itemize}[nosep]
  \item \emph{Rein:} Strommesswerte als CSV (vom Arduino oder aus einer
        Datei).
  \item \emph{Was passiert:} Jeder Messwert wird an \url{twin_core.py}
        übergeben.
  \item \emph{Stand:} mit Testdaten geprüft; Sensor noch nicht am
        Versuchsstand eingebaut.
\end{itemize}
\endgroup

%%%%%%%%%%%%%%%%%%%%%%%%%%%%%%%%%%%%%%%%%%%%%%%%%%%%%%%%%%%%%%%%%%%%%%%%%%%%%%%%
\section{Zusatzrechnungen}
\label{app:extra}

\anh{C.1 Vergleich numerischer Verfahren.} Finite Differenzen brauchen an
jeder Materialgrenze eine Sonderbehandlung; Randelemente führen zu voll
besetzten Matrizen und sind bei vielen Materialbereichen aufwendig; Finite
Volumen wären möglich. Die FEM behandelt Materialsprünge, Drehsymmetrie
($2\pi r$) und Konvektionsränder ohne Zusatzaufwand und ist Standard für
TEAM-Benchmarks.

\anh{C.2 Eigenwertanalyse.} Ohne Quelle gilt
$\mathbf M\dot{\mathbf T} = -\mathbf K\mathbf T$. Die Eigenformen aus
$\mathbf K\boldsymbol\varphi = \lambda\mathbf M\boldsymbol\varphi$ klingen
mit $\tau_i = 1/\lambda_i$ ab. Ergebnis: $\tau_1 = 204$\,s,
$\tau_2 = 4{,}4$\,s, Verhältnis 47 (Abbildung~\ref{fig:eigen}). Alle Formen
außer der ersten sind nach wenigen Sekunden abgeklungen; deshalb genügt eine
Form.

\begin{figure}[!ht]
\centering
\begin{tikzpicture}
\begin{axis}[width=\columnwidth, height=4.2cm, ybar, bar width=9pt,
    ymode=log, log origin=infty, ymin=0.1, ymax=1000,
    xtick={1,...,6}, xlabel={Eigenform $i$}, ylabel={$\tau_i$ [s]},
    ymajorgrids, grid style={gray!25}, font=\footnotesize,
    nodes near coords={\pgfmathprintnumber[precision=3, use comma]{\pgfplotspointmeta}\,s},
    point meta=rawy, every node near coord/.append style={font=\tiny},
    enlarge x limits=0.12]
  \addplot[fill=blue!60!cyan, draw=none] coordinates {(1,203.931)};
  \addplot[fill=gray!50, draw=none] coordinates
    {(2,4.38) (3,1.33) (4,0.637) (5,0.373) (6,0.246)};
\end{axis}
\end{tikzpicture}
\caption{Zeitkonstanten der sechs langsamsten Eigenformen\src{R5}.}
\label{fig:eigen}
\end{figure}

\anh{C.3 Reduziertes Modell gegen volle FEM.} Beide heizen die Scheibe bei
$5\,\mathrm A$ auf (Abbildung~\ref{fig:romfem}). Mit $\tau_1$ aus C.2 beträgt
der größte Fehler $0{,}36\,\mathrm K$ ($0{,}8\,\%$). Mit der im Code
verwendeten Zeitkonstante $\tau = C/UA = 244$\,s sind es bis zu
$3{,}0\,\mathrm K$ während des Aufheizens, beim Endwert~0. Grund: $UA$ wird
gegen $20\,^\circ$C gebildet, die Scheibenunterseite sieht aber von den
Spulen auf $30\,^\circ$C vorgewärmte Luft\src{R4, R6}.

\begin{figure}[!ht]
\centering
\begin{tikzpicture}
\begin{axis}[name=top, width=\columnwidth, height=4.6cm, xmin=0, xmax=20,
    ymin=0, xticklabels={}, ylabel={$\Delta T_{\max}$ [K]},
    grid=major, grid style={gray!25}, font=\footnotesize,
    legend pos=south east, legend style={font=\tiny, draw=none, fill=white, fill opacity=0.8, text opacity=1},
    legend cell align=left]
  \addplot[black, line width=1.6pt] coordinates {(0,0) (0.017,0.316) (0.083,1.332) (0.167,2.413) (0.250,3.423) (0.333,4.395) (0.417,5.339) (0.500,6.258) (0.583,7.154) (0.667,8.029) (0.750,8.882) (0.833,9.714) (0.917,10.527) (1.000,11.320) (1.083,12.093) (1.167,12.848) (1.250,13.585) (1.333,14.304) (1.417,15.006) (1.500,15.690) (2.000,19.464) (2.500,22.723) (3.000,25.537) (3.500,27.967) (4.000,30.066) (4.500,31.878) (5.000,33.442) (5.500,34.794) (6.000,35.960) (6.500,36.968) (7.000,37.838) (7.500,38.589) (8.000,39.237) (8.500,39.798) (9.000,40.281) (9.500,40.699) (10.000,41.060) (10.500,41.371) (11.000,41.640) (11.500,41.872) (12.000,42.073) (12.500,42.246) (13.000,42.395) (13.500,42.524) (14.000,42.636) (14.500,42.732) (15.000,42.815) (15.500,42.887) (16.000,42.949) (16.500,43.003) (17.000,43.049) (17.500,43.089) (18.000,43.123) (18.500,43.153) (19.000,43.179) (19.500,43.201) (20.000,43.220)};
  \addlegendentry{FEM, 697 Unbekannte (Referenz)}
  \addplot[blue!70!cyan, thick, dashed] coordinates {(0,0) (0.017,0.211) (0.083,1.047) (0.167,2.069) (0.250,3.066) (0.333,4.039) (0.417,4.989) (0.500,5.916) (0.583,6.820) (0.667,7.702) (0.750,8.563) (0.833,9.404) (0.917,10.224) (1.000,11.024) (1.083,11.805) (1.167,12.567) (1.250,13.310) (1.333,14.036) (1.417,14.744) (1.500,15.435) (2.000,19.244) (2.500,22.533) (3.000,25.373) (3.500,27.825) (4.000,29.943) (4.500,31.772) (5.000,33.351) (5.500,34.715) (6.000,35.892) (6.500,36.909) (7.000,37.787) (7.500,38.545) (8.000,39.200) (8.500,39.765) (9.000,40.253) (9.500,40.675) (10.000,41.038) (10.500,41.353) (11.000,41.624) (11.500,41.859) (12.000,42.061) (12.500,42.236) (13.000,42.387) (13.500,42.517) (14.000,42.629) (14.500,42.727) (15.000,42.810) (15.500,42.883) (16.000,42.945) (16.500,43.000) (17.000,43.046) (17.500,43.086) (18.000,43.121) (18.500,43.151) (19.000,43.177) (19.500,43.200) (20.000,43.219)};
  \addlegendentry{ROM, 1 Mode, $\tau_1$ aus Eigenproblem}
  \addplot[orange!90!red, thick, dash dot] coordinates {(0,0) (0.017,0.177) (0.083,0.878) (0.167,1.738) (0.250,2.581) (0.333,3.407) (0.417,4.216) (0.500,5.008) (0.583,5.785) (0.667,6.545) (0.750,7.291) (0.833,8.021) (0.917,8.736) (1.000,9.437) (1.083,10.124) (1.167,10.797) (1.250,11.456) (1.333,12.102) (1.417,12.735) (1.500,13.355) (2.000,16.820) (2.500,19.884) (3.000,22.595) (3.500,24.992) (4.000,27.112) (4.500,28.988) (5.000,30.646) (5.500,32.113) (6.000,33.411) (6.500,34.558) (7.000,35.573) (7.500,36.471) (8.000,37.264) (8.500,37.967) (9.000,38.588) (9.500,39.137) (10.000,39.623) (10.500,40.052) (11.000,40.432) (11.500,40.769) (12.000,41.066) (12.500,41.329) (13.000,41.561) (13.500,41.767) (14.000,41.949) (14.500,42.110) (15.000,42.252) (15.500,42.378) (16.000,42.489) (16.500,42.588) (17.000,42.675) (17.500,42.752) (18.000,42.820) (18.500,42.880) (19.000,42.933) (19.500,42.981) (20.000,43.022)};
  \addlegendentry{ROM im Code, $\tau=C/UA$}
\end{axis}
\begin{axis}[at={(top.south west)}, anchor=north west, yshift=-2mm,
    width=\columnwidth, height=3.0cm, xmin=0, xmax=20, ymin=0, ymax=3.5,
    xlabel={Zeit [min]}, ylabel={max.\ Fehler [K]},
    grid=major, grid style={gray!25}, font=\footnotesize]
  \addplot[blue!70!cyan, thick, dashed] coordinates {(0,0) (0.017,0.120) (0.083,0.296) (0.167,0.347) (0.250,0.358) (0.283,0.359) (0.333,0.357) (0.417,0.351) (0.500,0.344) (0.583,0.336) (0.667,0.328) (0.750,0.320) (0.833,0.312) (0.917,0.305) (1.000,0.297) (1.083,0.290) (1.167,0.283) (1.250,0.276) (1.333,0.270) (1.417,0.263) (1.500,0.257) (2.000,0.222) (2.500,0.191) (3.000,0.165) (3.500,0.143) (4.000,0.123) (4.500,0.106) (5.000,0.092) (5.500,0.079) (6.000,0.069) (6.500,0.059) (7.000,0.051) (7.500,0.044) (8.000,0.038) (8.500,0.033) (9.000,0.028) (9.500,0.025) (10.000,0.021) (10.500,0.018) (11.000,0.016) (11.500,0.014) (12.000,0.012) (12.500,0.010) (13.000,0.009) (13.500,0.008) (14.000,0.007) (14.500,0.006) (15.000,0.005) (15.500,0.004) (16.000,0.004) (16.500,0.003) (17.000,0.003) (17.500,0.002) (18.000,0.002) (18.500,0.002) (19.000,0.002) (19.500,0.001) (20.000,0.001)};
  \addplot[orange!90!red, thick, dash dot] coordinates {(0,0) (0.017,0.154) (0.083,0.465) (0.167,0.678) (0.250,0.844) (0.333,0.990) (0.417,1.125) (0.500,1.251) (0.583,1.371) (0.667,1.485) (0.750,1.592) (0.833,1.695) (0.917,1.792) (1.000,1.883) (1.083,1.970) (1.167,2.052) (1.250,2.130) (1.333,2.203) (1.417,2.272) (1.500,2.336) (2.000,2.645) (2.500,2.840) (3.000,2.943) (3.500,2.976) (3.533,2.976) (4.000,2.954) (4.500,2.890) (5.000,2.796) (5.500,2.681) (6.000,2.550) (6.500,2.410) (7.000,2.265) (7.500,2.118) (8.000,1.973) (8.500,1.831) (9.000,1.694) (9.500,1.562) (10.000,1.437) (10.500,1.319) (11.000,1.207) (11.500,1.104) (12.000,1.007) (12.500,0.917) (13.000,0.834) (13.500,0.757) (14.000,0.687) (14.500,0.622) (15.000,0.563) (15.500,0.509) (16.000,0.460) (16.500,0.415) (17.000,0.374) (17.500,0.337) (18.000,0.303) (18.500,0.273) (19.000,0.245) (19.500,0.220) (20.000,0.198)};
  \node[font=\tiny, anchor=west] at (axis cs:4.3,2.95) {3,0\,K};
  \node[font=\tiny, anchor=south west] at (axis cs:0.5,0.38) {0,36\,K};
\end{axis}
\end{tikzpicture}
\caption{Aufheizen der Scheibe bei $5\,\mathrm A$: Übertemperatur des
heißesten Knotens, volle FEM gegen ROM (oben), größter Fehler
im ganzen Feld (unten)\src{R6}.}
\label{fig:romfem}
\end{figure}

\anh{C.4 Magnetische Sättigung.} Bei konstantem $\mu_r$ ist die Flussdichte
proportional zum Strom. Linear hochgerechnet aus der FEM: Spitzenwert im
Eisen ca.\ $0{,}96\,\mathrm T$ bei $5\,\mathrm A$ eff.\ ($64\,\%$ von
$1{,}5\,\mathrm T$) und ca.\ $1{,}5\,\mathrm T$ bei $7{,}8\,\mathrm A$ eff.
Bei $7{,}8\,\mathrm A$ ist Reduktion A daher nur näherungsweise gültig. Die
lineare Hochrechnung überschätzt leicht\src{E5}.

\anh{C.5 Stromkonvention.} Gemessen wird der Effektivwert ($5\,\mathrm A$).
Die Verlustrechnung setzt ihn als Amplitude ein; die Wattzahlen in
Tabelle~\ref{tab:losses} sind daher um den Faktor 2 zu klein. Für die
Temperatur ist das unschädlich, weil der Faktor bei der Kalibrierung von
$hA$ aufgefangen wird. Die Kraftrechnung verwendet die echte Amplitude
$\sqrt2\cdot5 = 7{,}07\,\mathrm A$.

\anh{C.6 Skintiefe und Biot-Zahl.}
$\delta = \sqrt{2/(\omega\mu_0\sigma)} = 12{,}2$\,mm $\gg$ 3\,mm.
$\mathrm{Bi} = h\,(d/2)/k = 25\cdot0{,}0015/237 = 1{,}6\cdot10^{-4} \ll 0{,}1$:
Die Scheibe ist nahezu gleichmäßig warm.

\anh{C.7 Zeitkonstante.} $C = \rho c_p \pi R^2 d = 146{,}6$\,J/K;
$UA = P_{\mathrm{ref}}/\overline{\Delta T} = 25{,}81/43{,}03 = 0{,}600$\,W/K;
$\tau = C/UA = 244$\,s.

%%%%%%%%%%%%%%%%%%%%%%%%%%%%%%%%%%%%%%%%%%%%%%%%%%%%%%%%%%%%%%%%%%%%%%%%%%%%%%%%
\section{FEM-Details}
\label{app:fem}

\anh{Elemente.} Lineare Dreiecke. Mit Fläche $A_e$, Schwerpunktradius $r_c$
und konstanten Gradienten $b_k$, $c_k$ der Ansatzfunktionen:
\begin{equation*}
\mathbf K_e = 2\pi k\,(\mathbf b\mathbf b^\top + \mathbf c\mathbf c^\top)\,A_e\,r_c .
\end{equation*}

\anh{Konvektionsrand.} Randkante der Länge $L$ zwischen den Radien $r_a$,
$r_b$:
\begin{equation*}
\mathbf K_{\mathrm{Kante}} = 2\pi h\,\frac{L}{12}
\begin{pmatrix}3r_a+r_b & r_a+r_b\\ r_a+r_b & r_a+3r_b\end{pmatrix}.
\end{equation*}

\anh{Randbedingungen.} Feld: $A_\varphi = 0$ auf der Achse und am Außenrand
(Kontrolle mit Neumann-Rand). Wärme: Die Achse braucht keine Bedingung, da
$2\pi r$ dort verschwindet.

\anh{Löser.} Direkte LU-Zerlegung (\url{spsolve}); das Feldsystem wird
komplex gelöst. Die Eisenpermeabilität wird iterativ an die Flussdichte
angepasst (Konvergenz nach 2 Iterationen).

%%%%%%%%%%%%%%%%%%%%%%%%%%%%%%%%%%%%%%%%%%%%%%%%%%%%%%%%%%%%%%%%%%%%%%%%%%%%%%%%
\section{Reproduktion}
\label{app:repro}

Im Repository-Stamm:
{\footnotesize
\begin{verbatim}
python RUN.py            # rechnen + HTML-Ansicht
.venv/bin/python projectseminar_DT4TM_report/\
    report_numbers.py    # Rechenwerte (A.2)
.venv/bin/python projectseminar_DT4TM_report/\
    report_numbers.py --figure   # Abb. C1, C2
python refit_hA.py       # Kalibrierung hA
\end{verbatim}
}

%%%%%%%%%%%%%%%%%%%%%%%%%%%%%%%%%%%%%%%%%%%%%%%%%%%%%%%%%%%%%%%%%%%%%%%%%%%%%%%%
\section{Die Reduktion anschaulich}
\label{app:visual}

Abbildung~\ref{fig:reduction_visual} zeigt ohne Formeln, was die drei
Reduktionen aus Abschnitt~\ref{sec:mor} bewirken: Aus rund 9\,400 Zahlen
der FEM werden 8 Zahlen im ROM. Die drei Reduktionen sind \textbf{keine
Alternativen}: Sie wirken \textbf{gleichzeitig}, jede auf einen anderen Teil
des Geräts, und werden in jedem Zeitschritt nacheinander ausgewertet
(unterer Teil der Abbildung). Links steht jeweils das große Modell, in der
Mitte die Idee, rechts das reduzierte Modell.

\begin{figure*}[p]
\centering
\begin{tikzpicture}[font=\footnotesize, >={Stealth[length=2mm]},
    hdr/.style={font=\small\bfseries, text=white, fill=#1, rounded corners=2pt,
                minimum height=6mm, minimum width=5.4cm, inner sep=2pt},
    idee/.style={draw=gray!60, fill=yellow!8, rounded corners=3pt, align=left,
                 text width=5.0cm, inner sep=5pt},
    capt/.style={align=center, text width=5.2cm, font=\scriptsize},
    nd/.style={circle, draw, fill=white, minimum size=6.5mm, inner sep=0pt, font=\scriptsize},
    fl/.style={draw, rounded corners=3pt, align=center, fill=white, inner sep=3pt,
               minimum height=9mm},
    garr/.style={->, thick, gray!70}]
  % ---------- Kopfzeile ----------
  \node[hdr=gray!70!black] at (2.7,-0.4) {VORHER: großes Modell};
  \node[hdr=orange!80!black] at (8.7,-0.4) {IDEE};
  \node[hdr=green!45!black] at (14.7,-0.4) {NACHHER: ROM};

  % ================= Zeile A =================
  \foreach \i in {0,...,11}{\foreach \j in {0,...,4}{
    \pgfmathtruncatemacro{\c}{max(6,75-17*sqrt((\i-3)^2+(\j-1)^2))}
    \fill[orange!\c!white] (0.4+\i*0.383,-1.8-\j*0.4) rectangle ++(0.383,-0.4);
    \draw[gray!55, very thin] (0.4+\i*0.383,-1.8-\j*0.4) rectangle ++(0.383,-0.4);
    \draw[gray!55, very thin] (0.4+\i*0.383,-1.8-\j*0.4) -- ++(0.383,-0.4);
  }}
  \node[capt] at (2.7,-4.4) {Feld-FEM: \textbf{8\,700 komplexe Unbekannte}\\
    für \emph{jeden} Strom neu zu lösen};
  \node[idee] (IA) at (8.7,-3.0) {\textbf{Reduktion A: Strom-Skalierung}\\[2pt]
    Die Feldgleichung ist linear: doppelter Strom $\to$ vierfache Verluste,
    an jedem Ort gleich.\\[2pt]
    \emph{Wie ein Lautstärkeregler: dieselbe Melodie, nur lauter.}};
  \shade[left color=orange!80, right color=yellow!15] (12.1,-2.4) rectangle (13.5,-3.6);
  \node[font=\scriptsize] at (12.8,-2.15) {$\hat q(r,z)$ bei 5\,A};
  \node at (13.85,-3.0) {\large$\times$};
  \node[draw, rounded corners=2pt, fill=white] at (14.8,-3.0) {$(I/5\,\mathrm A)^2$};
  \node at (15.75,-3.0) {\large$=$};
  \shade[left color=orange!80, right color=yellow!15] (16.05,-2.4) rectangle (17.3,-3.6);
  \node[font=\scriptsize] at (16.7,-2.15) {$q(I)$};
  \node[capt] at (14.7,-4.4) {Verlustkarte \textbf{einmal} gespeichert,\\
    danach pro Strom nur \textbf{1 Multiplikation}};
  \draw[garr] (5.2,-3.0) -- (6.05,-3.0); \draw[garr] (11.35,-3.0) -- (12.0,-3.0);
  \draw[dashed, gray!40] (0,-5.05) -- (17.4,-5.05);

  % ================= Zeile B =================
  \foreach \k in {0,...,45}{
    \pgfmathtruncatemacro{\c}{30+40*exp(-((\k-28)/14)^2)}
    \fill[orange!\c!white] (0.4+\k*0.1,-7.2) rectangle ++(0.1,-0.6);
    \draw[gray!55, very thin] (0.4+\k*0.1,-7.2) rectangle ++(0.1,-0.6);
  }
  \node[font=\scriptsize] (t1) at (0.45,-6.55) {$T_1$};
  \node[font=\scriptsize] (t2) at (0.95,-6.55) {$T_2$};
  \node[font=\scriptsize] at (2.7,-6.55) {$\dots$};
  \node[font=\scriptsize] (tn) at (4.7,-6.55) {$T_{697}$};
  \draw[->, gray] (t1) -- (0.45,-7.15); \draw[->, gray] (t2) -- (0.55,-7.15);
  \draw[->, gray] (tn) -- (4.95,-7.15);
  \node[capt] at (2.7,-8.55) {Scheibe: \textbf{697 Knotentemperaturen},\\
    jede mit eigener Gleichung};
  \node[idee] (IB) at (8.7,-7.5) {\textbf{Reduktion B: eine Form für die Scheibe}\\[2pt]
    Die Scheibe ist fast überall gleich warm: Es ändert sich nur die
    \emph{Höhe} der Temperatur, nicht ihre \emph{Form}.\\[2pt]
    \emph{Wie ein Foto, das nur heller oder dunkler wird.}};
  \shade[left color=orange!30, right color=orange!30, middle color=orange!70] (12.1,-7.2) rectangle (15.1,-7.8);
  \node[font=\scriptsize] at (13.6,-6.9) {Form $\Delta T_{\mathrm{ref}}(r,z)$, einmal gespeichert};
  \node at (15.45,-7.5) {\large$\times$};
  \node[draw, rounded corners=2pt, fill=white, font=\normalsize] at (16.45,-7.5) {$\beta(t)$};
  \draw[thick] (12.1,-8.35) -- (15.1,-8.35);
  \draw (12.1,-8.25) -- (12.1,-8.45) node[below, font=\tiny] {$\beta=0$: kalt};
  \draw (15.1,-8.25) -- (15.1,-8.45) node[below, font=\tiny] {$\beta=1$: wie bei 5\,A};
  \fill[red!70!black] (13.9,-8.35) circle (1.3pt);
  \node[capt] at (14.7,-9.25) {\textbf{2 Zahlen}: $\beta_{\mathrm{Wirbel}}$ (Wirbelströme),
    $\beta_{\mathrm{Luft}}$ (warme Luft von unten)};
  \draw[garr] (5.2,-7.5) -- (6.05,-7.5); \draw[garr] (11.35,-7.5) -- (12.0,-7.5);
  \draw[dashed, gray!40] (0,-9.75) -- (17.4,-9.75);

  % ================= Zeile C =================
  \fill[gray!35] (0.4,-10.6) rectangle (1.49,-12.8);
  \fill[orange!35] (1.57,-10.6) rectangle (3.0,-12.8);
  \fill[gray!35] (3.13,-10.6) rectangle (3.76,-12.8);
  \fill[orange!35] (3.88,-10.6) rectangle (4.72,-12.8);
  \fill[cyan!12] (0.4,-10.1) rectangle (5.0,-10.6);
  \draw[step=0.18cm, gray!55, very thin] (0.4,-10.1) grid (5.0,-12.8);
  \node[capt] at (2.7,-13.3) {Spulen, Eisen, Luft als feines Feld:\\
    \textbf{Tausende Knoten} (so \emph{nicht} gerechnet)};
  \node[idee] (IC) at (8.7,-11.7) {\textbf{Reduktion C: RC-Netzwerk}\\[2pt]
    Jeder Körper ist fast gleich warm $\to$ ein Knoten pro Körper
    (Spule: Oberfläche + Wicklung).\\[2pt]
    \emph{Wie ein Schaltplan statt einer Landkarte.}};
  \node[nd, fill=orange!25] (iW) at (12.4,-11.1) {iW};
  \node[nd, fill=orange!25] (iO) at (13.55,-11.1) {iO};
  \node[nd, fill=gray!30]   (Fe) at (14.7,-11.1) {Fe};
  \node[nd, fill=orange!25] (aO) at (15.85,-11.1) {aO};
  \node[nd, fill=orange!25] (aW) at (17.0,-11.1) {aW};
  \node[nd, fill=cyan!12]   (Lu) at (14.7,-12.35) {L};
  \draw[<->] (iW) -- (iO); \draw[<->] (iO) -- (Fe); \draw[<->] (aO) -- (aW);
  \draw[->] (iO) -- (Lu); \draw[->] (Fe) -- (Lu); \draw[->] (aO) -- (Lu);
  \node[capt] at (14.7,-13.3) {\textbf{6 Knoten}, je eine Temperatur};
  \draw[garr] (5.2,-11.7) -- (6.05,-11.7); \draw[garr] (11.35,-11.7) -- (12.0,-11.7);

  % ================= Zusammenspiel =================
  \node[hdr=blue!55!black, minimum width=17.4cm] at (8.7,-14.3)
    {PRO ZEITSCHRITT: alle drei Reduktionen gleichzeitig};
  \node[fl] (I) at (1.3,-15.9) {Strom\\$I(t)$};
  \node[fl, fill=yellow!10] (A) at (4.6,-15.9) {\textbf{A}: Verluste\\$\times(I/5\,\mathrm A)^2$};
  \node[fl, fill=orange!10, text width=3.3cm] (C) at (9.2,-15.2) {\textbf{C}: 6 Knoten\\Spulen, Eisen, Luft};
  \node[fl, fill=orange!10, text width=3.3cm] (B) at (9.2,-16.7) {\textbf{B}: 2 Gewichte $\beta$\\Scheibe};
  \node[fl, fill=green!10] (O) at (14.3,-15.9) {Temperaturen\\+ Schwebehöhe};
  \draw[->, thick] (I) -- (A);
  \draw[->, thick] (A.east) -- ++(0.5,0) |- (C.west);
  \draw[->, thick] (A.east) -- ++(0.5,0) |- (B.west);
  \draw[->, thick, red!70!black] (C.south) -- node[left, font=\scriptsize, text=red!70!black] {warme Luft heizt Scheibe} (B.north);
  \draw[->, thick] (C.east) -- ++(0.6,0) |- (O.west);
  \draw[->, thick] (B.east) -- ++(0.6,0) |- (O.west);
  \node[draw=green!45!black, rounded corners=3pt, fill=green!5, align=center,
        text width=16.6cm, inner sep=4pt] at (8.7,-18.15)
    {Zustand nach jedem Schritt:\quad
     $x(t)=(\beta_{\mathrm{Wirbel}},\,\beta_{\mathrm{Luft}},\,T_{iO},\,T_{iW},\,T_{aO},\,T_{aW},\,T_{\mathrm{Fe}},\,T_{\mathrm L})$
     \quad$=$\quad\textbf{8 Zahlen}\\[2pt]
     statt $8\,700+697$ Unbekannten der FEM\quad(ein Zeitschritt $\approx 3\,\mu$s)};
\end{tikzpicture}
\caption{Die drei Reduktionen anschaulich. Oben: jede Reduktion von groß
(links) nach klein (rechts). Unten: Zusammenspiel in jedem Zeitschritt
(\url{twin_core.py}). Das feine Feld in Zeile C ist nur zur Veranschaulichung
gezeichnet; für Spulen und Eisen wurde keine Wärme-FEM gerechnet, das
RC-Netz wird direkt aus Geometrie und Werkstoffdaten aufgestellt.}
\label{fig:reduction_visual}
\end{figure*}

% Figure-Seite vor dem Literaturverzeichnis ausgeben.
\clearpage

\bibliographystyle{IEEEtran}
\bibliography{dt4tm_references}

\end{document}
